\documentclass[11pt]{article}

\usepackage{amsmath}
\usepackage{amssymb}
\usepackage{amsthm}
\usepackage{mathtools}
\usepackage{hyperref}

\theoremstyle{definition}
\newtheorem{definition}{Definition}
\newtheorem{example}{Example}

\theoremstyle{plain}
\newtheorem{theorem}{Theorem}
\newtheorem{proposition}[theorem]{Proposition}
\newtheorem{lemma}[theorem]{Lemma}

\title{Norms and Metrics}
\author{math-foundations / concept 09}
\date{}

\begin{document}
\maketitle

\section{Introduction}

A vector space gives us addition and scalar multiplication, but no notion
of \emph{length} or \emph{distance}. The moment we want to talk about
convergence of an iterative algorithm, continuity of a loss function, or
generalisation bounds on a hypothesis class, we need to measure how close
two vectors are. \emph{Norms} and \emph{metrics} provide exactly this
quantitative geometry.

The two notions are tightly related: every norm on a vector space induces
a metric in a canonical way. The converse fails, and the failure is
instructive --- it shows that metrics are strictly more general than
norms, and that some legitimate measures of distance (e.g.\ the discrete
metric, or the Kullback--Leibler divergence after suitable repair) cannot
be derived from any notion of vector length.

\section{Definitions}

Throughout, let $\mathbb{F}$ denote $\mathbb{R}$ or $\mathbb{C}$, and let
$V$ be a vector space over $\mathbb{F}$.

\begin{definition}[Norm]
A function $\|\cdot\| \colon V \to \mathbb{R}$ is a \emph{norm} if for
every $x,y \in V$ and every scalar $\alpha \in \mathbb{F}$,
\begin{enumerate}
  \item (Positive definiteness) $\|x\| \ge 0$, and
        $\|x\| = 0 \iff x = 0$.
  \item (Absolute homogeneity) $\|\alpha x\| = |\alpha|\,\|x\|$.
  \item (Triangle inequality) $\|x + y\| \le \|x\| + \|y\|$.
\end{enumerate}
The pair $(V, \|\cdot\|)$ is called a \emph{normed vector space}.
\end{definition}

\begin{definition}[Metric]
Let $X$ be a set. A function $d \colon X \times X \to \mathbb{R}$ is a
\emph{metric} (or \emph{distance function}) if for every $x,y,z \in X$,
\begin{enumerate}
  \item (Non-negativity) $d(x,y) \ge 0$.
  \item (Identity of indiscernibles) $d(x,y) = 0 \iff x = y$.
  \item (Symmetry) $d(x,y) = d(y,x)$.
  \item (Triangle inequality) $d(x,z) \le d(x,y) + d(y,z)$.
\end{enumerate}
The pair $(X, d)$ is called a \emph{metric space}.
\end{definition}

\begin{definition}[Induced metric]
Let $(V, \|\cdot\|)$ be a normed vector space. The \emph{metric induced
by} $\|\cdot\|$ is
\[
  d(x,y) := \|x - y\|, \qquad x,y \in V.
\]
\end{definition}

\begin{example}[$\ell^p$ norms on $\mathbb{R}^n$]
For $1 \le p < \infty$ and $x = (x_1,\dots,x_n) \in \mathbb{R}^n$, define
\[
  \|x\|_p := \left( \sum_{i=1}^n |x_i|^p \right)^{1/p},
  \qquad
  \|x\|_\infty := \max_{1 \le i \le n} |x_i|.
\]
Each $\|\cdot\|_p$ is a norm. The triangle inequality for $p \ge 1$ is the
\emph{Minkowski inequality}.
\end{example}

\section{Main theorem}

\begin{theorem}[Norms induce metrics]
\label{thm:norm-induces-metric}
Let $(V, \|\cdot\|)$ be a normed vector space. Then the function
\[
  d \colon V \times V \to \mathbb{R}, \qquad d(x,y) := \|x - y\|,
\]
is a metric on $V$.
\end{theorem}

\begin{proof}
We verify the four axioms of a metric.

\smallskip
\noindent\emph{(1) Non-negativity.}
By positive definiteness of the norm,
$d(x,y) = \|x - y\| \ge 0$ for all $x,y \in V$.

\smallskip
\noindent\emph{(2) Identity of indiscernibles.}
We use positive definiteness again:
\[
  d(x,y) = 0
  \iff \|x - y\| = 0
  \iff x - y = 0
  \iff x = y.
\]

\smallskip
\noindent\emph{(3) Symmetry.}
By absolute homogeneity with $\alpha = -1$,
\[
  d(x,y)
  = \|x - y\|
  = \bigl\|(-1)\,(y - x)\bigr\|
  = |-1|\,\|y - x\|
  = \|y - x\|
  = d(y,x).
\]

\smallskip
\noindent\emph{(4) Triangle inequality.}
For any $x,y,z \in V$, write $x - z = (x - y) + (y - z)$. The triangle
inequality for the norm gives
\[
  d(x,z)
  = \|x - z\|
  = \bigl\|(x - y) + (y - z)\bigr\|
  \le \|x - y\| + \|y - z\|
  = d(x,y) + d(y,z).
\]

All four axioms hold, so $d$ is a metric on $V$.
\end{proof}

\section{A metric that is not norm-induced}

\begin{proposition}[Discrete metric]
\label{prop:discrete}
Let $V$ be a nonzero vector space and define
\[
  d(x,y) =
  \begin{cases}
    0 & \text{if } x = y, \\
    1 & \text{if } x \ne y.
  \end{cases}
\]
Then $d$ is a metric on $V$, but there is no norm $\|\cdot\|$ on $V$
such that $d(x,y) = \|x - y\|$ for all $x,y$.
\end{proposition}

\begin{proof}
That $d$ is a metric is a direct check of the four axioms.

Suppose for contradiction that $d(x,y) = \|x - y\|$ for some norm
$\|\cdot\|$. Pick any $x \ne 0$. By homogeneity,
\[
  \|2x\| = 2\,\|x\|.
\]
On the other hand, $2x \ne 0$ and $x \ne 0$, so
\[
  \|2x\| = d(2x, 0) = 1 \qquad \text{and} \qquad
  \|x\|  = d(x, 0)  = 1.
\]
Combining the displays gives $1 = 2$, a contradiction.
Hence $d$ does not come from any norm.
\end{proof}

\section{Why this matters}

Once a metric is fixed, the entire machinery of analysis is unlocked:
open and closed sets, continuity, convergence, Cauchy sequences,
completeness, compactness, Lipschitz constants. In machine learning, the
choice of norm is rarely innocent: $\ell^1$ regularisation produces
sparse solutions, $\ell^2$ regularisation shrinks coefficients
isotropically, and the spectral (operator) norm controls how much a
neural-network layer can amplify inputs. The Wasserstein distance is a
genuine metric on probability distributions and underlies optimal
transport; the Kullback--Leibler divergence is \emph{not} symmetric and
\emph{not} a metric, even though it is often informally called a
``distance''.

\end{document}
